\section{Related work}

\paragraph{Decision interfaces and constrained outputs.}
TypeSafe introduces Jev as a System One model with typed decisions \citep{almeida2026jev}; our use of that name refers to the interface. Laya supplies the open implementation used in this study \citep{laya2026}. The biological question is whether a shared candidate scorer can combine reusable task inputs with reliable predictions, rather than whether label constraints alone are new. Constrained autoregressive decoding already restricts permitted outputs \citep{scholak2021picard}; exact likelihood scoring over a finite label list is another necessary control. We therefore distinguish unambiguous output identity from correct biological interpretation and include constrained generation in the prospective design.

\paragraph{Sequence--text alignment and instruction tuning.}
Joint use of biological sequences and natural language predates this study. ProtST augments protein representation learning with biomedical descriptions and evaluates supervised and zero-shot prediction \citep{xu2023protst}. BioTranslator maps biological features and textual descriptions into a shared space for zero-shot biomedical classification \citep{xu2023biotranslator}. These methods motivate explicit attention to label semantics, but our two fixed label sets do not test their zero-shot setting. A candidate string can be present in an input without the model demonstrating transferable semantic understanding.

Instruction-tuned biological models broaden the interface further. LLaMA-Gene combines vocabulary expansion, additional sequence pretraining, and instruction fine-tuning \citep{wang2024llamagene}. Its instruction corpus is also a source of the two BioPAWS-2 task views used here. ChatNT connects biological-sequence inputs with a conversational language interface \citep{almeida2025chatnt}. Our experiment evaluates a finite candidate distribution from an encoder rather than open-ended answer generation, and omits additional neural sequence pretraining. The differing data, objectives, and evaluation settings preclude cross-paper accuracy rankings.

\paragraph{Tokenization and encoder adaptation.}
DNABERT-2 uses BPE in a genome foundation model and evaluates its representation and computational benefits \citep{zhou2024dnabert}. That evidence does not imply that adding a biological vocabulary to a text-derived checkpoint will improve a small supervised adaptation. Our representation comparison holds examples and update counts fixed while recording the larger embedding matrix and measured training cost. Laya builds a typed-decision interface over an encoder; the checkpoint used here employs ModernBERT-large \citep{warner2025modernbert,laya2026}. We retain the checkpoint's 1,024-token configuration rather than infer a longer usable context from the underlying encoder family.

\paragraph{Calibration and evaluation scope.}
Temperature scaling adjusts confidence after training without changing the most probable class \citep{guo2017calibration}. We fit temperatures on a separate calibration split and report test probability metrics alongside classification metrics. Sequence interventions and candidate permutations address different questions: the former probe dependence under changed inputs, whereas the latter preserve the candidate set and allow semantic alignment of predictions. Neither substitutes for homology-aware partitioning or evaluation on new task descriptions.
